\documentclass[11pt]{article}
\usepackage[T1]{fontenc}
\usepackage{lmodern}
\usepackage[margin=1in]{geometry}
\usepackage{amsmath,amssymb,amsthm,mathtools}
\usepackage{enumitem}
\usepackage{booktabs,tabularx,array}
\usepackage[hidelinks]{hyperref}
\usepackage{microtype}

\newtheorem{theorem}{Theorem}
\newtheorem{proposition}[theorem]{Proposition}
\newtheorem{definition}[theorem]{Definition}
\newtheorem{remark}[theorem]{Remark}

\title{From the Accepted Package to an Independent Odd-$m$ D5 Theorem Package}
\author{}
\date{2026-03-14}

\begin{document}
\maketitle

\begin{abstract}
This note answers one packaging question and does not search for a new theorem.
The odd-$m$ $d=5$ globalization theorem is already closed inside the accepted
chain
\[
033 \to 062 \to 076 \to 077 \to 079 \to 081 \to 082 \to 083.
\]
No new mathematical bottleneck is visible in that chain. What remains is to
rewrite the imported inputs so that the final theorem no longer reads as proved
only ``within the accepted package.'' The main missing item is now one compact
concrete bridge theorem that replaces the still-imported \texttt{076} package
in manuscript order. A one-pass readable downstream globalization package is
now available in \texttt{099}, but it still leaves the standing theorem-package
inputs explicit at the top.
\end{abstract}

\section{Purpose and exact current status}

This is a cleanup note, not a new theorem note. Its job is to separate two
claims that can otherwise get blurred together:
\begin{enumerate}[label=\textup{(\arabic*)},leftmargin=2.2em]
\item the odd-$m$ $d=5$ theorem is already closed \emph{inside the accepted
package};
\item the theorem is not yet written as one independently self-contained
manuscript-order package.
\end{enumerate}

\begin{theorem}[Current accepted-package theorem]
For odd $m$ in the accepted D5 regime, every actual lift of the exceptional
cutoff row $\delta=3m-3$ continues through
\[
3m-2 \to 3m-1
\]
and lies in the regular continuing endpoint class there. Consequently:
\begin{enumerate}[label=\textup{(\arabic*)},leftmargin=2.2em]
\item there is no mixed-status realized $\delta$;
\item fixed realized $\delta$ determines tail length;
\item fixed realized $\delta$ determines endpoint class and future word;
\item $\rho=\rho(\delta)$ globally;
\item raw global $(\beta,\delta)$ is exact on the true accessible boundary
union.
\end{enumerate}
\end{theorem}

\begin{remark}[Scope reminder]
The theorem above is exactly the current accepted conclusion, but its proof is
still distributed across imported notes. In particular, \texttt{076} remains a
\emph{componentwise} concrete bridge theorem, and the accepted-chain proof
still passes through \texttt{062}, \texttt{077}, \texttt{079}, \texttt{081},
and \texttt{083}. In the cleaned package, the exact \texttt{079},
\texttt{081}, \texttt{077}, and separate \texttt{062} roles are now supplied
by the compact notes \texttt{095}--\texttt{098}, leaving \texttt{076} as the
main imported theorem layer.
\end{remark}

\section{What an independent package should mean}

\begin{definition}[Independent odd-$m$ D5 theorem package]
An \emph{independent package} is a manuscript-order theorem package in which
all proof-critical objects used by the final odd-$m$ $d=5$ globalization
argument are stated in place, with their exact scope visible, so that the
final theorem no longer depends on the phrase ``within the accepted package''
for its mathematical status.
\end{definition}

This notion of independence is deliberately modest. It does \emph{not} ask for
new numerics, a new search, or a new idea. It asks for one self-contained
presentation layer in which the proof-critical imports have become theorems of
the package itself.

\section{What already looks theorem-level enough}

\begin{proposition}[Stable layers]
The following pieces already look strong enough to survive essentially
unchanged in an independent package:
\begin{enumerate}[label=\textup{(\arabic*)},leftmargin=2.2em]
\item the abstract bridge $(\beta,\rho)$;
\item the componentwise concrete bridge $(\beta,\delta)$;
\item the reduction from fixed-$\delta$ ambiguity to full-chain tail length;
\item the regular-union theorem;
\item the exceptional-row reduction;
\item the final gluing proof pattern.
\end{enumerate}
\end{proposition}

So the missing work is not discovery. It is exposing the already-used inputs in
one clean order.

\section{What is still missing}

\subsection*{1. One self-contained structural theorem block}

\begin{proposition}[Needed repackaging: structural input]
An independent package should contain one theorem block that states, in one
place:
\begin{enumerate}[label=\textup{(\alph*)},leftmargin=2.2em]
\item the exact trigger family
\[
H_{L1}=\{(q,w,u,\lambda)=(m-1,m-1,u,2):u\neq 2\};
\]
\item the candidate active orbit on current coordinates;
\item the phase-$1$ invariant
\[
q\equiv u-\rho+1 \pmod m;
\]
\item the universal first-exit targets
\[
T_{\mathrm{reg}}=(m-1,m-2,1),
\qquad
T_{\mathrm{exc}}=(m-2,m-1,1);
\]
\item the conclusion that, on the actual active branch, the true orbit agrees
with the candidate orbit up to first exit.
\end{enumerate}
\end{proposition}

At present this content is mathematically sound but distributed across
\texttt{033} and \texttt{062}. Independence asks for one theorem block, not a
new proof idea.

\begin{remark}[Update after 098]
This structural block is now written explicitly in
\texttt{theorem/d5\_098\_compact\_cleanup\_033\_062\_structural\_block.md}.
The only remaining issue there is the broader defect-template provenance behind
the trigger-family lemma, not a missing structural theorem for the odd-$m$
globalization package.
\end{remark}

\subsection*{2. One explicit chart-to-raw continuation theorem}

\begin{proposition}[Needed repackaging: actual exceptional continuation]
An independent package should include one theorem that packages the current
composition of \texttt{062} and \texttt{079} into a single proof object:

\smallskip
\noindent every actual lift of the exceptional cutoff row reaches the universal
raw first-exit target
\[
T_{\mathrm{exc}}=(m-2,m-1,1)
\]
by direction $1$, and therefore continues in chain labels through
\[
3m-3\to 3m-2\to 3m-1,
\]
where $3m-2$ is the terminal regular source-$4$ occurrence and $3m-1$ is the
regular source-$1$ start.
\end{proposition}

\begin{remark}[Why this is the key packaging gap]
This is the one place where the current proof still feels assembled from
accepted notes rather than written as a single theorem object. This theorem is
weaker than the final endpoint-class gluing theorem: it packages only the
actual exceptional continuation, not the later tail-length or regular-union
arguments that promote that continuation to global $\rho=\rho(\delta)$.
\end{remark}

\begin{remark}[Update after 095 and 098]
The cleaned suite now has this role through the \texttt{095} chart/interface
reproof together with the explicit structural block in \texttt{098}. So this is
no longer the main missing item in the independent-package gap.
\end{remark}

\subsection*{3. One compact concrete bridge theorem}

\begin{proposition}[Needed repackaging: concrete bridge]
An independent package should contain one compact theorem that states, in one
place, the intrinsic boundary coordinates
\[
\sigma\equiv w+u-q-1\pmod m,
\qquad
\delta=q+m\sigma,
\]
the splice law
\[
\delta' \equiv \delta+1 \pmod{m^2},
\]
the current-event readout from $(\beta,\delta)$, and the component-image
statement that each splice-connected accessible component maps to one forward
$+1$ orbit segment in $\mathbb Z/m^2\mathbb Z$.
\end{proposition}

This content already exists in \texttt{076}. The missing work is only to
present it as one compact theorem in manuscript order.

\subsection*{4. One theorem-level definition block}

\begin{proposition}[Needed repackaging: final definitions]
An independent package should define in one place the final proof objects used
later without reconstruction from earlier notes:
\begin{enumerate}[label=\textup{(\alph*)},leftmargin=2.2em]
\item the accessible boundary union;
\item a splice-connected accessible component;
\item a full regular chain;
\item the exceptional cutoff row;
\item an actual lift;
\item endpoint class, padded future word, and right-congruence state.
\end{enumerate}
\end{proposition}

The notions are already present in the chain, but they are spread across
notes. An independent package should not force the reader to reconstruct them
from context.

\begin{remark}[Update after 092]
The cleaned theorem suite now effectively contains this definition block, so
this is no longer a serious remaining gap.
\end{remark}

\subsection*{5. Support files should become verification aids only}

\begin{proposition}[Needed repackaging: move proof-critical content out of
support files]
The support files
\begin{center}
\ttfamily
checks/d5\_077\_compact\_interval\_summary.json \quad and \quad
checks/d5\_081\_regular\_union\_endpoint\_table.csv
\end{center}
should remain attached as verification aids, but any proof-critical content now
navigated through them should be restated as lemmas or propositions in the
manuscript-order theorem notes.
\end{proposition}

Again, this is a packaging requirement, not evidence of a missing theorem idea.

\section{Recommended manuscript order}

A compact independent package can be organized in the following order:
\begin{enumerate}[label=\textup{\arabic*.},leftmargin=2.2em]
\item the cleaned definition block already present in \texttt{092};
\item the compact structural theorem note \texttt{098};
\item one concrete bridge theorem note replacing the current promoted-support
status of \texttt{076};
\item the compact cleanup notes \texttt{097}, \texttt{096}, and \texttt{095};
\item the final globalization theorem.
\end{enumerate}

That order makes the final proof read as a forward argument rather than as a
reconstruction from accepted imports.

\section{Smallest honest checklist}

The shortest path from ``closed inside the accepted package'' to ``proved in a
self-contained package'' is:
\begin{enumerate}[label=\textup{\arabic*.},leftmargin=2.2em]
\item keep \texttt{098} as the compact structural theorem block;
\item compress \texttt{076} into one final concrete bridge theorem note;
\item keep the \texttt{095}--\texttt{097} cleanup notes as the stable compact
support layers;
\item keep the JSON and CSV files as support only, not as implicit proof steps;
\item keep the broader \texttt{033} defect-template provenance explicit only
where it is genuinely needed.
\end{enumerate}

If those five items are done, the phrase
\begin{quote}
closed inside the accepted package
\end{quote}
should be replaceable by
\begin{quote}
proved in a self-contained odd-$m$ D5 theorem package.
\end{quote}

\section{Bottom line}

The remaining gap is no longer a theorem-search gap. It is a theorem-packaging
gap:
\begin{enumerate}[label=\textup{(\arabic*)},leftmargin=2.2em]
\item one compact concrete bridge theorem statement replacing \texttt{076}, and
\item optionally, a first-principles proof of the broader \texttt{033}
trigger-family classification beyond what \texttt{098} isolates for the
globalization chain.
\end{enumerate}

So the right next work is theorem cleanup, not new theorem search.

\end{document}
